\documentclass[a4paper,11pt]{article}
\usepackage[a4paper,margin=0.68in,top=0.63in,bottom=0.60in]{geometry}
% Compile with XeLaTeX (Overleaf: Menu > Compiler > XeLaTeX).
\usepackage{fontspec}
\usepackage{xeCJK}
\setmainfont[Scale=0.97]{Liberation Sans}
\setsansfont[Scale=0.97]{Liberation Sans}
\setCJKmainfont{Noto Sans CJK TC}
\setCJKsansfont{Noto Sans CJK TC}
\renewcommand{\familydefault}{\sfdefault}
\usepackage{xcolor,tabularx,enumitem,titlesec,microtype,hyperref,fancyhdr}
\hypersetup{hidelinks}

\definecolor{ink}{HTML}{182333}
\definecolor{muted}{HTML}{536173}
\definecolor{rulecol}{HTML}{CBD3DF}
\definecolor{shade}{HTML}{F5F7FA}
\color{ink}
\setlength{\parindent}{0pt}
\setlength{\parskip}{0pt}
\linespread{1.045}
\setlength{\tabcolsep}{0pt}

\pagestyle{fancy}
\fancyhf{}
\fancyfoot[L]{\footnotesize\color{muted} Yu-Jui Liang（梁栯睿）}
\fancyfoot[R]{\footnotesize\color{muted} \thepage\ / 2}
\renewcommand{\headrulewidth}{0pt}
\renewcommand{\footrulewidth}{0pt}
\setlength{\footskip}{17pt}

\newcommand{\dochead}[2]{%
  {\footnotesize\bfseries\color{muted} RESEARCH \& PROJECT SUMMARY}%
  \hfill{\footnotesize\color{muted} #1}\par\vspace{4pt}
  {\fontsize{20}{23}\selectfont\bfseries Yu-Jui Liang\hspace{0.4em}{\fontsize{17}{20}\selectfont 梁栯睿}}\par\vspace{3pt}
  {\small\color{muted}#2}\par\vspace{7pt}
  \color{rulecol}\rule{\linewidth}{1.1pt}\color{ink}\vspace{10pt}}
\newcommand{\smallhead}[1]{%
  {\footnotesize\bfseries\color{muted} #1}\par\vspace{4pt}
  \color{rulecol}\rule{\linewidth}{0.8pt}\color{ink}\vspace{10pt}}
\newcommand{\metadata}[1]{{\small #1}\par\vspace{9pt}}
\newcommand{\project}[6]{%
  {\small\color{muted}\bfseries #1}\par\vspace{3pt}
  {\fontsize{12}{14}\selectfont\bfseries #2}\par\vspace{7pt}
  {\small\textbf{Purpose.}\ #3}\par\vspace{5pt}
  {\small\textbf{Approach.}\ #4}\par\vspace{5pt}
  {\small\textbf{Evidence / Status.}\ #5}\par\vspace{5pt}
  {\footnotesize\textbf{Relevant skills:}\ #6}\par\vspace{8pt}}
\newcommand{\flowline}[1]{%
  \noindent\colorbox{shade}{%
    \parbox{\dimexpr\linewidth-2\fboxsep}{\centering\footnotesize\color{muted} #1}}
  \par\vspace{12pt}}
\newcommand{\separate}{%
  \par\noindent\color{rulecol}\rule{\linewidth}{0.7pt}\color{ink}\par\vspace{11pt}}

% FIGURE POLICY (not rendered in this two-page HR summary):
%   1. MO mapping: original technical figure showing single-output vs.
%      multi-output implementation, or a VERIFIED QoR comparison plot.
%   2. STA ECO: gate sizing / VT swapping / buffer insertion
%      candidate search, branch list, and global-best rollback (methodology).
%   3. Optional for physical-design roles: actual post-route Mini-SoC image
%      with a readable caption. Do not add a diffusion image without a
%      meaningful comparison or verified experimental result.
% Prefer a separate one-page visual appendix over shrinking this handout.
\begin{document}
\dochead{EDA / CAD / DIGITAL IC}{Logic Synthesis \quad / \quad Physical Design \quad / \quad CAD Flow Automation}
% Education / advisor metadata (Chinese labels; tailored EDA / CAD JD keywords remain English).
\metadata{%
學歷：國立臺灣大學｜資料科學碩士學位學程（預計 2027 年 6 月畢業）\\[3pt]
\phantom{學歷：}國立政治大學｜資訊科學系學士（雙主修：電子物理學士學位學程）\\[3pt]
指導教授：江介宏教授（Jie-Hong Roland Jiang），國立臺灣大學\\[3pt]
共同指導：鄭湘筠博士（Hsiang-Yun Cheng），中央研究院\\[3pt]
專業領域：EDA / CAD Flow Development, Logic Synthesis / Technology Mapping, RTL-to-GDS,\\[3pt]
\phantom{專業領域：}Physical Design / APR, Floorplanning, STA / Timing ECO, PPA / QoR Optimization,\\[3pt]
\phantom{專業領域：}SoC Integration, Tcl / Python Automation, ML for EDA}

% ===========================================================================
% PROJECT 01 / 04 -- LOGIC SYNTHESIS & TECHNOLOGY MAPPING
% ABC-based MO mapping: custom MO cut discovery, Boolean matching,
% area-aware MILP pair selection, then differentiable mapping optimization.
% QoR: 6% average ADP improvement vs. single-output mapping, ASAP7 7 nm.
% ===========================================================================
\project{PROJECT 01 \quad / \quad LOGIC SYNTHESIS \& TECHNOLOGY MAPPING}
{Multi-Output Standard-Cell Mapping with MILP and Differentiable Optimization}
{Improve area--delay product (ADP) by exploiting multi-output standard cells in the academic ASAP7 7-nm technology library, compared with conventional single-output mapping.}
{Extended the ABC logic-synthesis tool in C++ with a custom multi-output (MO) cut-discovery path and Boolean matching. Applied area-aware mixed-integer linear programming (MILP) to select matched cut/cell pairs, then passed the selected candidates to a PyTorch-based differentiable mapping optimizer for Liberty-based area/delay evaluation.}
{Validated correctness and evaluated 30 arithmetic benchmarks. The updated MILP-enhanced mapping flow achieved a 6\% average ADP improvement over the single-output baseline.}
{C++, ABC, logic synthesis, technology mapping, Boolean matching, PyTorch, Liberty, ASAP7, QoR analysis}
\flowline{ABC / MO Cut Discovery $\longrightarrow$ Boolean Matching $\longrightarrow$ Area-Aware MILP Selection $\longrightarrow$ Differentiable Mapping}
\separate

% ===========================================================================
% PROJECT 02 / 04 -- STA-DRIVEN INCREMENTAL TIMING ECO
% Primary ECO operations: critical-path gate sizing, VT swapping, and
% selective buffer insertion. Relocation is not the focus of this project.
% Prior-flow evidence: five public benchmarks; 16.7% median setup TNS,
% 14.6% average setup WNS, ~8% less power, ~2% greater HPWL.
% These verified numbers PRECEDE the newer bounded candidate-search policy.
% ===========================================================================
\project{PROJECT 02 \quad / \quad TIMING ECO \& CAD AUTOMATION}
{STA-Driven Incremental Timing ECO Flow (ICCAD 2025 CAD Contest Problem C)}
{Improve setup timing through STA-driven gate sizing, threshold-voltage (VT) swapping, and targeted buffer insertion, while balancing timing improvements against power and physical-design constraints.}
{Developed an OpenROAD-based ECO loop that analyzes critical-path STA reports to generate gate-sizing, VT-swapping, and buffer-insertion candidates. Each trial undergoes incremental physical updates and fresh parasitic-aware STA, with setup WNS/TNS, power, HPWL, and placement legality rechecked. A bounded best-first search explores alternative ECO sequences through a branch list, while a global-best checkpoint retains the best observed QoR for rollback.}
{Evaluated five public benchmarks, achieving 16.7\% median
setup TNS improvement and 14.6\% average setup WNS
improvement relative to seed DEFs. Power consumption
decreased by approximately 8\%, with a modest 2\%
increase in HPWL.}
{OpenROAD, Tcl/Python scripting, STA, Physical Design, Timing ECO, PPA analysis}
\flowline{Critical-Path STA $\longrightarrow$ Sizing / VT Swap / Buffering $\longrightarrow$ Candidate Search $\longrightarrow$ PPA Validation}

\newpage
\smallhead{SELECTED PROJECTS \quad | \quad EDA / CAD / DIGITAL IC}

% ===========================================================================
% PROJECT 03 / 04 -- RTL-TO-GDS & PHYSICAL DESIGN
% Reusable end-to-end integration baseline for new EDA tools/scripts/methods,
% complementing isolated benchmarks; keep independently verified PPA facts.
% ===========================================================================
\project{PROJECT 03 \quad / \quad RTL-TO-GDS \& PHYSICAL DESIGN}
{Open-Source Mini-SoC and Macro-Based Physical Implementation}
{Build a reusable RTL-to-GDS baseline to validate that new EDA tools, scripts, and methodologies run correctly in an end-to-end physical-design flow, beyond standalone benchmark tests.}
{Integrated a Wishbone-based Mini-SoC with SPI/I2C peripherals, an SRAM controller, and a 1-KiB SRAM hard macro; implemented the design in ASAP7 and extended the flow to an RV32I system with four 1-KiB SRAM macros. Used these designs as integration testbeds for tool interoperability and script execution across design stages.}
{The Mini-SoC contained 11.1k instances at a 500-ps clock; post-route results achieved +34.8 ps setup slack with zero setup/hold violations. The RV32I design integrated four 1-KiB SRAM macros.}
{Verilog, SoC integration, SRAM macro integration, APR, STA, timing closure, ASAP7}
\flowline{RTL / SRAM Integration $\longrightarrow$ RTL-to-GDS Baseline $\longrightarrow$ Tool / Script Integration Checks}
\separate

% ===========================================================================
% PROJECT 04 / 04 -- ML FOR EDA
% FloorSet-Lite provides 1M TRAINING layouts with near-optimal wirelength;
% avoid claiming mathematical/global optimality for all training examples.
% No verified aggregate QoR gain; report the method rather than a result.
% MT-tree terminology not verified; retain documented B*-tree representation.
% ===========================================================================
\project{PROJECT 04 \quad / \quad ML FOR EDA}
{Diffusion-Guided SoC Floorplanner (ICCAD 2026 CAD Contest Problem C)}
{Improve macro-floorplanning QoR by reducing sensitivity to initial placement quality, which can strongly affect the final HPWL and constraint satisfaction of local search.}
{Used a diffusion model pretrained on one million near-optimal FloorSet-Lite floorplans to generate higher-quality initial macro placements. Fed the generated placements into constraint-aware B*-tree simulated annealing for HPWL/area refinement and macro legalization.}
{Developed a two-stage diffusion-initialization and annealing-refinement pipeline for the FloorSet challenge, targeting feasible macro placements and better floorplan QoR; no quantified improvement is claimed here.}
{PyTorch, ML for EDA, floorplanning, macro placement, B*-tree, simulated annealing, HPWL, legalization}
\flowline{Pretrained Diffusion (1M Samples) $\longrightarrow$ B*-tree / Annealing Refinement $\longrightarrow$ Legalization / QoR}

\vfill
\end{document}
